\cleardoublepage
\chapter*{Provas Iconográficas --- Handouts}
\addcontentsline{toc}{chapter}{Provas Iconográficas --- Handouts}
\markboth{Provas Iconográficas --- Handouts}{}

\begin{rpgbox}[ARQUIVO DE MESA]
\small
As imagens deste caderno pertencem à ficção documental do romance. Não são reproduções de documentos históricos nem provas de arquivo. Foram compostas para este volume; o horror real de Barbacena e da seca não é reconstituído aqui.
\end{rpgbox}

\clearpage
\newcommand{\prova}[4]{%
\begin{center}
\begin{rpgcard}[PROVA #1 · #2]
\centering
\includegraphics[width=0.62\textwidth,height=0.50\textheight,keepaspectratio]{figures/immersao/#3}\\[0.4em]
{\footnotesize\itshape #4}
\vspace{0.25em}
{\scriptsize\scshape\color{rpgblood} COTA · ARQ. 1.260 · PROVA #1}
\end{rpgcard}
\end{center}
\vfill\clearpage}
\prova{01}{PACIENTE 1.260}{placa_00.png}{Paciente 1.260. Retrato de arquivo.\\Reconstituição ficcional.}
\prova{02}{ADMISSÃO}{placa_01.png}{Registro de admissão. Colônia de Barbacena, 1917.\\Reconstituição ficcional.}
\prova{03}{O MÉDICO}{placa_02.png}{Dr. Heitor Oliveira. Arquivo pessoal, 1979.\\Reconstituição ficcional.}
\prova{04}{O HOMEM}{placa_03.png}{O homem do pano preto. Imagem não catalogada.\\Reconstituição ficcional.}
\prova{05}{A INVESTIGADORA}{placa_04.png}{Dra. Lúcia Mendes. Rodoviária de Barbacena, 2026.\\Reconstituição ficcional.}
\prova{06}{A FOTOGRAFIA}{placa_05.png}{Fotografia 5. Arquivo pessoal de Lúcia Mendes, 2026.\\Reconstituição ficcional.}
